% META: {"id":"1SPE-DERLOCAL-CO-003","chapitre":"1SPE-DERIVATION-LOCAL","type_objet":"corrige","exercice_id":"1SPE-DERLOCAL-EX-003","status":"approved","origin":"generated","status_history":["generated","audited","accepted","approved"],"release_acceptance":"RELEASE_OWNER_BATCH_ACCEPTANCE","acceptance_closure_digest":"sha256:9b3ccf9a81c5520fbb7e03b7b2d3e7bf2057a02834b6d908bf3be5f49f3b553f","acceptance_receipt_digest":"sha256:4fad86f0282e0fa4e0419cca1ad6379ae7af5a280c04a4a90880f90a1c044242"}
% BEGIN-VERIFY
% from sympy import Rational
% assert Rational(80-0, 2-0) == 40
% assert Rational(275-80, 5-2) == 65
% END-VERIFY
\begin{corrige}{1SPE-DERLOCAL-EX-003}
\textbf{1.} On calcule le taux de variation de $d$ entre $0$ et $2$ :
\[
\frac{d(2)-d(0)}{2-0}=\frac{80-0}{2}=40.
\]
Ce résultat signifie que le véhicule a parcouru en moyenne $40$~km par heure sur les deux premières heures. L'unité est le kilomètre par heure (km/h).

\textbf{2.} Entre $2$ et $5$ :
\[
\frac{d(5)-d(2)}{5-2}=\frac{275-80}{3}=65.
\]
Sur les trois heures suivantes, la vitesse moyenne est de $65$~km/h.

\textbf{3.} On compare les deux taux de variation : $65>40$. Le véhicule a roulé plus vite en moyenne sur l'intervalle $[2\,;\,5]$ que sur $[0\,;\,2]$.

\textit{Vérification :} les deux taux sont positifs, ce qui est cohérent car la distance ne peut que croître.
\end{corrige}
